% ═══════════════════════════════════════════════════════════════════════
%  Section 6 — Security Measures
% ═══════════════════════════════════════════════════════════════════════

\section{Security Measures and Threat Mitigations}
\label{sec:security}

This section details the specific security mechanisms implemented and the threats they mitigate.

\subsection{Authentication Security}

\begin{table}[H]
\centering
\begin{tabular}{p{4cm} p{4cm} p{5cm}}
    \toprule
    \textbf{Threat} & \textbf{Mitigation} & \textbf{Implementation} \\
    \midrule
    Password cracking & Argon2id (memory-hard) & 64~MB memory, 3 iterations, 4 threads \\
    Credential stuffing & Account lockout & 5 attempts → 15 min lock \\
    User enumeration & Timing-safe responses & Dummy hash on unknown emails \\
    Token theft (XSS) & HttpOnly cookies & Tokens never in JS-accessible storage \\
    Session hijacking & Device fingerprinting & IP + User-Agent hash bound to token \\
    Token replay & Refresh rotation & New JTI on each refresh; reuse → full revocation \\
    TOTP replay & Used-OTP table & Code + time window tracked in DB \\
    \bottomrule
\end{tabular}
\caption{Authentication threat model and mitigations.}
\end{table}

\subsection{Web Application Security}

\subsubsection{Cross-Site Scripting (XSS)}

\begin{itemize}
    \item \textbf{Server-side:} All user text input is sanitized with \texttt{bleach.clean()} before storage. The sanitizer strips all HTML tags, JavaScript URIs, and event handlers.
    \item \textbf{Client-side:} React's JSX automatically escapes variables, preventing DOM-based XSS.
    \item \textbf{Headers:} \texttt{X-Content-Type-Options: nosniff} prevents MIME-type sniffing attacks.
\end{itemize}

\subsubsection{Cross-Site Request Forgery (CSRF)}

The platform implements the \textbf{double-submit cookie pattern}:
\begin{enumerate}
    \item \texttt{GET /api/auth/csrf} returns a random token in both a cookie and the response body.
    \item The frontend reads the body value and sends it as the \texttt{X-CSRF-Token} header on state-changing requests.
    \item The backend middleware verifies that the header value matches the cookie value.
    \item Safe methods (GET, HEAD, OPTIONS) are exempt.
\end{enumerate}

\subsubsection{SQL Injection}

\begin{itemize}
    \item All database queries use \textbf{SQLAlchemy ORM} with parameterized queries.
    \item No raw SQL strings are used anywhere in the codebase.
    \item Pydantic schemas validate and type-check all inputs before they reach the ORM layer.
\end{itemize}

\subsubsection{Session Fixation}

\begin{itemize}
    \item Sessions are created server-side with fresh UUIDs — the client cannot influence session IDs.
    \item On login, new session and tokens are always created (no session reuse).
    \item On logout, the session is revoked in the database and cookies are cleared.
\end{itemize}

\subsection{Data Protection}

\subsubsection{Encryption at Rest}

\begin{table}[H]
\centering
\begin{tabular}{p{3.5cm} p{3.5cm} p{6cm}}
    \toprule
    \textbf{Data} & \textbf{Method} & \textbf{Details} \\
    \midrule
    Passwords & Argon2id hash & One-way; never stored as plaintext \\
    Resumes (PDF) & Fernet (AES-128-CBC) & Encrypted on upload, decrypted in-memory on download \\
    TOTP secrets & Fernet & Encrypted before DB storage; decrypted for verification only \\
    \bottomrule
\end{tabular}
\caption{Data encryption methods.}
\end{table}

\subsubsection{File Upload Security}

Resume and avatar uploads are protected by a multi-layer validation pipeline:

\begin{enumerate}
    \item \textbf{Extension whitelist:} Only \texttt{.pdf} for resumes; \texttt{.jpg}, \texttt{.jpeg}, \texttt{.png} for avatars.
    \item \textbf{MIME type check:} Content-Type header must match allowed types.
    \item \textbf{Magic byte validation:} File header bytes must match expected signatures (prevents extension spoofing).
    \item \textbf{Size limits:} 2~MB for both resumes and avatars.
    \item \textbf{Content scan (PDF):} Checks for embedded JavaScript, launch actions, and auto-open actions.
    \item \textbf{UUID filenames:} Original filenames are stored only in the database; files on disk use random UUIDs (prevents path traversal).
    \item \textbf{Storage isolation:} Files are stored outside the web root; access requires authentication and authorization.
\end{enumerate}

\subsection{Security Headers}

The following HTTP security headers are set on all responses:

\begin{lstlisting}
X-Content-Type-Options: nosniff
X-Frame-Options: DENY
X-XSS-Protection: 1; mode=block
Referrer-Policy: strict-origin-when-cross-origin
Strict-Transport-Security: max-age=31536000; includeSubDomains
\end{lstlisting}

\subsection{Access Control Model}

\begin{table}[H]
\centering
\begin{tabular}{p{3cm} p{2.5cm} p{2.5cm} p{2.5cm}}
    \toprule
    \textbf{Resource} & \textbf{User} & \textbf{Recruiter} & \textbf{Admin} \\
    \midrule
    Own profile & Read/Write & Read/Write & Read/Write \\
    Others' profiles & — & Read & Read \\
    Own resumes & CRUD & CRUD & CRUD \\
    Others' resumes & — & Read & Read \\
    Job postings & Read & CRUD & CRUD \\
    User management & — & — & Full \\
    Audit logs & — & — & Read \\
    \bottomrule
\end{tabular}
\caption{RBAC permission matrix.}
\end{table}
